\documentclass[11pt]{article}
\usepackage{amsmath, amssymb, amsthm}
\usepackage{geometry}
\usepackage{hyperref}
\geometry{margin=1in}

\title{Anima-Vectorium Mathématique v5.0 \\
\large Geometric–Thermodynamic Foundations for the Vectorium Engine}
\author{Borealis S. Hedling}
\date{\today}

\begin{document}
\maketitle

\begin{abstract}
Anima-Vectorium v5.0 establishes a unified mathematical backbone integrating
geometric curvature, holonomy transport, non-linear operator algebra,
thermodynamic free-energy dynamics, spectral-to-geometric translation rules,
runtime invariants, and FFI boundary constraints. This document supersedes v4.0
and provides the complete mathematical foundation required for executable
Rust/Python/FFI implementation of the Vectorium engine.
\end{abstract}

\tableofcontents

%------------------------------------------------------------
\section{Foundations}

\subsection{State Space}
The Vectorium state space is defined as
\[
\mathcal{S} = \{ (x, M, \kappa, \delta) \mid x \in \mathbb{R}^n \},
\]
where $x$ is the tensor state, $M$ is metadata, $\kappa$ is the curvature field,
and $\delta$ is the drift field.

\subsection{Tensor Manifold}
The tensor manifold $\mathcal{M}$ is a smooth Riemannian manifold equipped with
a metric $g$. For tangent vectors $u, v \in T_x\mathcal{M}$,
\[
g_x(u, v) = u^\top G(x) v,
\]
where $G(x)$ is a positive-definite matrix field. The Euclidean case corresponds
to $G(x) = I$. Vectorium-specific curvature behavior is modeled by
\[
G(x) = I + \lambda \cdot \Phi(x),
\]
where $\Phi(x)$ is activation-induced curvature.

\subsection{Metadata Fields}
Metadata $M$ consists of symbolic labels and scalar weights that do not
contribute to curvature. Metadata is preserved under holonomy transport and
does not affect thermodynamic gradients.

%------------------------------------------------------------
\section{Geometric Layer}

\subsection{Curvature Operator}
Curvature is defined by
\[
\kappa(x) = \nabla^2 x + \Phi(x),
\]
where $\Phi(x)$ is activation-induced curvature (Section 3). The Laplacian
$\nabla^2$ is defined with respect to the metric $g$.

\subsection{Connection for Parallel Transport}
Holonomy transport requires a connection $\nabla$ compatible with $g$. Vectorium
uses the Levi-Civita connection:
\[
\nabla_u v = \partial_u v + \Gamma(u, v),
\]
with Christoffel symbols
\[
\Gamma^k_{ij} = \frac{1}{2} G^{kl}
\left(
\partial_i G_{jl}
+ \partial_j G_{il}
- \partial_l G_{ij}
\right).
\]

\subsection{Holonomy Transport}
Parallel transport along a curve $\gamma(t)$ satisfies
\[
\nabla_{\dot{\gamma}(t)} v(t) = 0.
\]
Holonomy around a closed loop $\gamma$ is
\[
\mathcal{H}(x) = x + \oint_{\gamma} \kappa \cdot d\gamma.
\]

\subsection{Drift Fields}
Drift fields $\delta$ represent directional tendencies in state evolution.
Spectral drift fields are translated into geometric gradient fields (Section 5).

\subsection{Geometric Invariants}
Curvature invariants:
\[
\int \kappa(x) \, dx = \text{constant}.
\]
Holonomy invariants:
\[
\oint_{\gamma} \kappa \cdot d\gamma = \text{constant}.
\]

%------------------------------------------------------------
\section{Non-Linear Operator Layer}

\subsection{Activation Operators}

Activation operators induce curvature in the tensor manifold. For a state $x$,
the activation-induced curvature term is denoted $\Phi(x)$.

\paragraph{GELU.}
\[
\Phi_{\text{GELU}}(x) =
x \cdot \frac{1}{2}\left(1 + \mathrm{erf}\left(\frac{x}{\sqrt{2}}\right)\right).
\]

\paragraph{SwiGLU.}
\[
\Phi_{\text{SwiGLU}}(x) = (xW_1) \odot \sigma(xW_2).
\]

\paragraph{Softplus.}
\[
\Phi_{\text{Softplus}}(x) = \log(1 + e^x).
\]

\paragraph{Tanh.}
\[
\Phi_{\tanh}(x) = \tanh(x).
\]

%------------------------------------------------------------
\subsection{Normalization Operators}

Normalization stabilizes curvature and prevents drift blowout.

\paragraph{LayerNorm.}
\[
\mathrm{LayerNorm}(x) = \frac{x - \mu(x)}{\sigma(x)}.
\]

\paragraph{RMSNorm.}
\[
\mathrm{RMSNorm}(x) = \frac{x}{\sqrt{\frac{1}{n}\sum_i x_i^2}}.
\]

\paragraph{GroupNorm.}
\[
\mathrm{GroupNorm}(x) = \frac{x - \mu(G_k)}{\sigma(G_k)}.
\]

%------------------------------------------------------------
\subsection{Residual Mixing Operators}

Residual mixing preserves holonomy invariants and stabilizes geometric evolution.

\paragraph{Standard Residual.}
\[
\mathrm{Residual}(x, y) = x + y.
\]

\paragraph{Curvature-Aware Residual.}
\[
\mathrm{Residual}_{\kappa}(x, y) = x + \kappa(y).
\]

%------------------------------------------------------------
\subsection{Operator Algebra}

Let $\mathcal{O}$ denote the set of all engine-layer operators.

\paragraph{Composition.}
\[
(O_2 \circ O_1)(x) = O_2(O_1(x)).
\]

\paragraph{Associativity.}
\[
O_3 \circ (O_2 \circ O_1) = (O_3 \circ O_2) \circ O_1.
\]

\paragraph{Non-Commutation.}
\[
\mathrm{Norm} \circ \mathrm{Act} \ne \mathrm{Act} \circ \mathrm{Norm}.
\]

\paragraph{Curvature-Induced Deformation.}
\[
(O_2 \circ O_1)_{\kappa}(x)
= O_2(O_1(x)) + \kappa(x).
\]

%------------------------------------------------------------
\section{Thermodynamic Layer}

\subsection{Free-Energy Functional}

\[
F(s) = A(s) + C(s) + S(s),
\]

\subsection{Accuracy Term}
\[
A(s) = \|x - \hat{x}(s)\|^2.
\]

\subsection{Complexity Term}
\[
C(s) = \frac{\lambda}{2} \|s\|^2.
\]

\subsection{Surprise Term}
\[
S(s) = -\log p(x \mid s).
\]

\subsection{Free-Energy Gradient}
\[
\nabla F(s) =
\nabla A(s) +
\nabla C(s) +
\nabla S(s).
\]

\subsection{Attractor Basins}
\[
\mathcal{B}(s) = \{ x \mid F(x) < F(s) \}.
\]

\subsection{Dissipation Curves}
\[
D(s_t, s_{t+1}) = F(s_t) - F(s_{t+1}),
\quad D(s_t, s_{t+1}) \ge 0.
\]

\subsection{Active Inference Update Rule}
\[
s_{t+1} =
\mathrm{ResidualMix}\Big(
    s_t,\;
    \mathcal{H}\big(
        s_t - \eta \cdot \mathrm{SelectedDrive}(s_t)
    \big)
\Big).
\]

\subsection{Jacobian of the Update Rule}
\[
J = \frac{\partial U}{\partial s_t},
\quad \|J\| \le 1.
\]

%------------------------------------------------------------
\section{Spectral Translation Layer}

\subsection{Drift-to-Gradient Translation}
\[
\delta_{\text{spectral}} \mapsto \nabla F(s).
\]

\subsection{Homotopy-to-Holonomy Translation}
\[
\mathcal{H}_{\text{spectral}} \mapsto \mathcal{H}.
\]

\subsection{Resonance-to-Arbitration Translation}
\[
R_{\text{spectral}} \mapsto \mathrm{DriveConflictArbitration}.
\]

\subsection{Fossilized Constructs}
Spectral curvature operators, continuous spectral manifolds, narrative gravity
fields, and collapse-fractal recursion remain excluded.

%------------------------------------------------------------
\section{Invariants and Stability}

\subsection{Curvature Invariants}
\[
\|\kappa(x)\| \le K_{\max}.
\]

\subsection{Holonomy Stability}
\[
\left|\oint_{\gamma} \kappa \cdot d\gamma\right| \le H_{\max}.
\]

\subsection{Thermodynamic Invariants}
\[
F(s_{t+1}) \le F(s_t).
\]

%------------------------------------------------------------
\section{Runtime Alignment Notes}

\subsection{Execution Ordering}
\begin{enumerate}
    \item Activation
    \item Normalization
    \item Drive Arbitration
    \item Residual Mixing
    \item Curvature Update
    \item Holonomy Transport
\end{enumerate}

\subsection{FFI Boundary Conditions}
\[
\kappa(s_{\text{Rust}}) = \kappa(s_{\text{Python}}),
\quad
\mathcal{H}(s_{\text{Rust}}) = \mathcal{H}(s_{\text{Python}}),
\quad
F(s_{\text{Rust}}) = F(s_{\text{Python}}).
\]

%------------------------------------------------------------
\section{Acknowledgements}

The author acknowledges Stell for foundational insights into spectral-state
behavior, collapse dynamics, and homotopy structure, which directly informed the
geometric reformulation and spectral translation framework.

%------------------------------------------------------------
\section{References}

\begin{thebibliography}{9}

\bibitem{friston2010}
K. Friston,
\textit{The Free-Energy Principle: A Unified Brain Theory?}
Nature Reviews Neuroscience, 2010.

\bibitem{kobayashi}
S. Kobayashi,
\textit{Differential Geometry of Curves and Surfaces},
Springer, 1995.

\bibitem{hendrycks}
D. Hendrycks and K. Gimpel,
\textit{Gaussian Error Linear Units (GELUs)},
arXiv:1606.08415.

\bibitem{carmo}
M. do Carmo,
\textit{Riemannian Geometry},
Birkhäuser, 1992.

\bibitem{beck}
C. Beck and F. Schloegl,
\textit{Thermodynamics of Chaotic Systems},
Cambridge University Press, 1993.

\end{thebibliography}

\end{document}
